\documentclass[11pt]{article}

\usepackage{amsmath,amssymb}
\usepackage{xurl}
\usepackage[hidelinks]{hyperref}
\setlength{\emergencystretch}{2em}

% Shared commands for notes in docs/paper-gaps/.

\DeclareUnicodeCharacter{2080}{\ifmmode{}_{0}\else\textsubscript{0}\fi}
\DeclareUnicodeCharacter{2081}{\ifmmode{}_{1}\else\textsubscript{1}\fi}
\DeclareUnicodeCharacter{2082}{\ifmmode{}_{2}\else\textsubscript{2}\fi}
\DeclareUnicodeCharacter{2099}{\ifmmode{}_{n}\else\textsubscript{n}\fi}

\newcommand{\Lean}{\textsc{Lean}}
\ExplSyntaxOn
\NewDocumentCommand{\leanid}{m}{
  \ifmmode
    \textsf{\detokenize{#1}}
  \else
    \group_begin:
    \urlstyle{sf}
    \tl_set:Nn \l_tmpa_tl {#1}
    \tl_replace_all:Nnn \l_tmpa_tl {\_} {_}
    \exp_args:NV \path \l_tmpa_tl
    \group_end:
  \fi
}
\ExplSyntaxOff
\providecommand{\tr}{\operatorname{tr}}
\newcommand{\tnleanIssueBase}{https://github.com/LionSR/TNLean/issues/}
\newcommand{\tnleanPrBase}{https://github.com/LionSR/TNLean/pull/}
\newcommand{\tnleanDocBase}{https://sirui-lu.com/TNLean/docs/}
\newcommand{\ghissue}[1]{\href{\tnleanIssueBase#1}{GitHub issue \##1}}
\newcommand{\ghpr}[1]{\href{\tnleanPrBase#1}{PR \##1}}
\newcommand{\doclink}[1]{\href{\tnleanDocBase#1.html}{\path{#1}}}

% Dirac notation and the identity operator, as in blueprint/src/macros/common.tex.
% \providecommand keeps note-local definitions authoritative.
\usepackage{dsfont}
\providecommand{\ket}[1]{|#1\rangle}
\providecommand{\bra}[1]{\langle #1|}
\providecommand{\braket}[2]{\langle #1 | #2 \rangle}
\providecommand{\ketbra}[2]{|#1\rangle\!\langle #2|}
\providecommand{\Id}{\mathds{1}}
\providecommand{\one}{\mathds{1}}
\providecommand{\C}{\mathbb{C}}
\providecommand{\MN}[1]{M_{#1}(\C)}
\providecommand{\MD}{M_D(\C)}

% External referencing for paper sources.  The build helper writes sanitized
% auxiliary files under build/paper-aux/ and excludes duplicated source labels.
\usepackage{xr}

% Import a generated paper auxiliary file with a caller-supplied label prefix.
% The candidate paths support builds from docs/paper-gaps/, the repository root,
% and build/paper-aux/ itself.
\newcommand{\importpaperaux}[2]{%
  \IfFileExists{../../build/paper-aux/#2.aux}{%
    \externaldocument[#1]{../../build/paper-aux/#2}%
  }{%
    \IfFileExists{build/paper-aux/#2.aux}{%
      \externaldocument[#1]{build/paper-aux/#2}%
    }{%
      \IfFileExists{#2.aux}{%
        \externaldocument[#1]{#2}%
      }{}%
    }%
  }%
}

% General external reference macros
\newcommand{\paperref}[2]{\ref{#1-#2}}
\newcommand{\papereqref}[2]{(\ref{#1-#2})}

% CPSV16 source (1606.00608 / MPDO-22-12-17-2.tex) external document and shortcuts
\importpaperaux{cpsv16-}{1606.00608/MPDO-22-12-17-2-tnlean}
\newcommand{\cpsvref}[1]{\paperref{cpsv16}{#1}}
\newcommand{\cpsveqref}[1]{\papereqref{cpsv16}{#1}}

% Verdict marker: \gapnote{<kind>}{<status>}. Kind is the policy
% classification (clarification, local-correction, scope-restriction,
% unfaithful, false-source, open-gap); status is open, wip (elimination
% underway), resolved, or historical. Machine-read by the site index;
% prints a verdict line.
\newcommand{\gapnote}[2]{\par\noindent\textbf{Verdict:} #1 (#2).\par}


\title{Strong Subadditivity as a Standalone Axiom and Its Lieb-Concavity
Elimination Route}
\author{}
\date{2026-06-22}

\begin{document}
\maketitle
\gapnote{open-gap}{resolved}

\section{The assertion and its formal status}

Let $\rho_{ABC}$ be a density operator on the finite-dimensional Hilbert space
$\mathcal H_A\otimes\mathcal H_B\otimes\mathcal H_C$, and let
$\rho_{AB}$, $\rho_{BC}$, and $\rho_B$ denote its reduced density operators.
The von Neumann entropy is
\begin{align}
    S(\rho)
        &= -\tr(\rho\log\rho).
    \label{eq:cpsv16_ssa_from_lieb_route_entropy_definition}
\end{align}
Strong subadditivity states
\begin{align}
    S(\rho_{ABC})+S(\rho_B)
        &\leq S(\rho_{AB})+S(\rho_{BC}).
    \label{eq:cpsv16_ssa_from_lieb_route_strong_subadditivity}
\end{align}
Lieb and Ruskai proved this inequality in finite dimensions
\cite{LiebRuskai1973}.

The formal development previously asserted
Eq.~\ref{eq:cpsv16_ssa_from_lieb_route_strong_subadditivity} as the standalone
axiom \leanid{strong_subadditivity}. It now derives the result as
\leanid{strong_subadditivity_general} in
\doclink{TNLean/Channel/Schwarz/StrongSubadditivityPosDef}. The former axiom has
been removed from \doclink{TNLean/Entropy/SSAEqualityCharacterization}, and its
single re-export, \leanid{Entropy.strongSubadditivity}, now points to the
derived theorem.\footnote{The formal entropy in
\doclink{TNLean/Analysis/Entropy} is defined from the eigenvalues by
$S(\rho)=\sum_i\eta(\lambda_i)$, where $\eta(x)=-x\log x$.}

The derivation uses Lieb's joint concavity theorem
\cite{Lieb1973Convex}.\footnote{The formal input is
\leanid{lieb_concavity_posDef} in
\doclink{TNLean/Analysis/LiebConcavity}
(\path{TNLean/Analysis/LiebConcavity.lean}). It is consumed through
\leanid{lieb_concavity} in
\doclink{TNLean/Channel/Schwarz/AndoLieb}.}
The only entropy-side dependency of
\leanid{strong_subadditivity_general} that replaces the former axiom is the
proved theorem \leanid{lieb_concavity_posDef}. This note records the complete
route from that theorem to strong subadditivity.

\section{The mathematical route}

The derivation passes through quantum relative entropy and its monotonicity
under a partial trace. Each layer below states the required mathematical fact
and records its formal implementation.

\begin{enumerate}
    \item \emph{Entropy bridge.}

    The eigenvalue definition of entropy agrees with
    Eq.~\ref{eq:cpsv16_ssa_from_lieb_route_entropy_definition}. This bridge
    permits the trace and matrix-logarithm calculations used below.\footnote{The
    entropy bridge is tracked on \ghissue{784}.}

    \item \emph{Quantum relative entropy.}

    For density operators $\rho$ and $\sigma$ satisfying
    $\ker\sigma\subseteq\ker\rho$, define
    \begin{align}
        D(\rho\,\| \sigma)
            &= \tr(\rho(\log\rho-\log\sigma)).
        \label{eq:cpsv16_ssa_from_lieb_route_relative_entropy_definition}
    \end{align}
    The support condition ensures that the logarithm of $\sigma$ is evaluated
    only where $\rho$ has support. This definition is formalized as
    \leanid{quantumRelativeEntropy}.

    \item \emph{Klein's inequality.}

    Quantum relative entropy satisfies
    \begin{align}
        D(\rho\,\| \sigma)
            &\geq 0,
        \label{eq:cpsv16_ssa_from_lieb_route_klein_nonnegative}\\
        D(\rho\,\| \sigma)=0
            &\Longleftrightarrow \rho=\sigma.
        \label{eq:cpsv16_ssa_from_lieb_route_klein_equality}
    \end{align}
    The nonnegativity statement on the full support domain is
    \leanid{quantumRelativeEntropy_nonneg_general} in
    \doclink{TNLean/Analysis/KleinInequality}
    (\path{TNLean/Analysis/KleinInequality.lean}). Its full-rank base case,
    \leanid{quantumRelativeEntropy_nonneg}, reduces the problem in eigenbases
    to a doubly stochastic overlap matrix and the scalar Gibbs inequality
    $\log x\leq x-1$.

    For the singular case, the reference density operator is regularized as
    \begin{align}
        \sigma'_\varepsilon
            &= \frac{\sigma+\varepsilon\Id}{1+\varepsilon N},
        \qquad \varepsilon>0,
        \label{eq:cpsv16_ssa_from_lieb_route_reference_regularization}
    \end{align}
    where $N$ is the matrix dimension. The matrix
    $\sigma'_\varepsilon$ is positive definite and has the same eigenvectors
    as $\sigma$. As $\varepsilon\to0^+$, the matrix limit reduces to scalar
    logarithm limits on the positive eigenvalues. The condition
    $\ker\sigma\subseteq\ker\rho$ eliminates the diagonal weights associated
    with zero eigenvalues. The formal limit is
    \leanid{tendsto_re_trace_mul_log_perturb}. This proof uses neither Lieb
    concavity nor operator Jensen.

    The full-rank equality case is
    \leanid{quantumRelativeEntropy_eq_zero_iff}. If the relative entropy
    vanishes, every term in the Gibbs slack is tight. The strict inequality
    $\log x<x-1$ for $x\neq1$ forces the eigenvalues to agree on the overlap
    support. The resulting relation makes the overlap matrix intertwine the
    eigenvalue diagonals, and conjugating $\sigma$ into the eigenbasis of
    $\rho$ gives $\rho=\sigma$. This argument also uses neither Lieb concavity
    nor operator Jensen.

    \item \emph{Joint convexity.}

    For $t\in[0,1]$, joint convexity means
    \begin{align}
        D\bigl(t\rho_0+(1-t)\rho_1
          \|\,t\sigma_0+(1-t)\sigma_1\bigr)
        &\leq tD(\rho_0\,\| \sigma_0)
          +(1-t)D(\rho_1\,\| \sigma_1).
        \label{eq:cpsv16_ssa_from_lieb_route_joint_convexity}
    \end{align}
    This is the only layer that consumes Lieb's joint concavity theorem
    \cite{Lieb1973Convex}. For $s\in[0,1)$, define the Lindblad--Uhlmann
    approximant
    \begin{align}
        g_s(\rho,\sigma)
            &= \frac{\tr\rho
              -\operatorname{Re}\tr(\rho^s\sigma^{1-s})}{1-s}.
        \label{eq:cpsv16_ssa_from_lieb_route_relative_entropy_approximant}
    \end{align}
    Lieb's concavity of
    $(A,B)\mapsto\tr(K^\dagger A^sKB^{1-s})$ implies that each $g_s$ is jointly
    convex. The approximants satisfy
    \begin{align}
        g_s(\rho,\sigma)
            &\longrightarrow D(\rho\,\| \sigma)
            \qquad\text{as }s\to1^-.
        \label{eq:cpsv16_ssa_from_lieb_route_approximant_limit}
    \end{align}
    Their pointwise limit is therefore convex.

    The positive-definite result is
    \leanid{convexOn_quantumRelativeEntropy} in
    \doclink{TNLean/Channel/Schwarz/RelativeEntropyConvexity}
    (\path{TNLean/Channel/Schwarz/RelativeEntropyConvexity.lean}).
    The convexity of the approximants is
    \leanid{convexOn_relativeEntropyApprox}, which consumes
    \leanid{lieb_concavity_posDef}, and their convergence is
    \leanid{tendsto_relativeEntropyApprox}. The double-sum identity
    \leanid{re_trace_cfc_mul_cfc_eq_double_sum} used in the limit requires no
    Lieb input.

    The theorem
    \leanid{convexOn_quantumRelativeEntropy_support} extends joint convexity to
    the full domain $\ker\sigma\subseteq\ker\rho$. It regularizes both
    arguments by
    \begin{align}
        M_\varepsilon
            &= \frac{M+\varepsilon\Id}{1+\varepsilon N},
        \qquad \varepsilon>0,
        \label{eq:cpsv16_ssa_from_lieb_route_affine_trace_one_perturbation}
    \end{align}
    applies positive-definite convexity to the four regularized endpoints and
    their mixture, and passes to the limit with
    \leanid{tendsto_relativeEntropyPerturb}. The domain is convex because the
    kernel of a positively weighted sum of positive semidefinite matrices is
    the intersection of their kernels; this step is
    \leanid{mulVec_eq_zero_of_smul_add}. The extension uses no additional Lieb
    input.

    \item \emph{Data processing under a partial trace.}

    For a trace-preserving completely positive map $\Lambda$, data processing
    is the inequality
    \begin{align}
        D(\Lambda\rho\,\| \Lambda\sigma)
            &\leq D(\rho\,\| \sigma).
        \label{eq:cpsv16_ssa_from_lieb_route_data_processing}
    \end{align}
    Only the case in which $\Lambda$ is a partial trace is required here. Its
    proof combines joint convexity with unitary invariance and
    ancilla-additivity:
    \begin{align}
        D(U\rho U^\dagger\,\|\,U\sigma U^\dagger)
            &= D(\rho\,\| \sigma),
        \label{eq:cpsv16_ssa_from_lieb_route_unitary_invariance}\\
        D(\rho\otimes\tau\,\| \sigma\otimes\tau)
            &= D(\rho\,\| \sigma),
        \label{eq:cpsv16_ssa_from_lieb_route_ancilla_additivity}
    \end{align}
    where $U$ is unitary and $\tau$ is a unit-trace ancilla state. Joint
    convexity and unitary invariance alone do not imply partial-trace
    monotonicity; Eq.~\ref{eq:cpsv16_ssa_from_lieb_route_ancilla_additivity} is
    also required.

    Unitary invariance is formalized for arbitrary Hermitian arguments as
    \leanid{quantumRelativeEntropy_conj_unitary} in
    \doclink{TNLean/Channel/Schwarz/RelativeEntropyUnitaryInvariance}
    (\path{TNLean/Channel/Schwarz/RelativeEntropyUnitaryInvariance.lean}).
    It follows from covariance of the matrix logarithm,
    \leanid{Matrix.log_conj_unitary}, which is the $f=\log$ case of
    \leanid{Matrix.cfc_conj_unitary}, and from trace cyclicity. This result
    requires no Lieb input.

    On positive-definite arguments, ancilla-additivity is
    \leanid{quantumRelativeEntropy_kronecker} in
    \doclink{TNLean/Channel/Schwarz/RelativeEntropyAncillaAdditivity}
    (\path{TNLean/Channel/Schwarz/RelativeEntropyAncillaAdditivity.lean}). Its
    main identity is
    \begin{align}
        \log(\rho\otimes\tau)
            &= \log\rho\otimes\Id+\Id\otimes\log\tau.
        \label{eq:cpsv16_ssa_from_lieb_route_tensor_logarithm}
    \end{align}
    This is \leanid{Matrix.log_kronecker}. The one-sided identities
    \leanid{Matrix.cfc_kronecker_one} and
    \leanid{Matrix.cfc_one_kronecker} follow from functional calculus through
    the unital tensor embeddings. Their recombination uses the logarithm of a
    commuting product, \leanid{Matrix.PosDef.cfc_log_mul}, in
    \doclink{TNLean/Analysis/CfcLogAdditive}
    (\path{TNLean/Analysis/CfcLogAdditive.lean}). The
    $\Id\otimes\log\tau$ terms cancel in the difference of tensor logarithms,
    the trace of a Kronecker product factors, and $\tr\tau=1$ gives
    Eq.~\ref{eq:cpsv16_ssa_from_lieb_route_ancilla_additivity}.

    The singular-support version is
    \leanid{quantumRelativeEntropy_kronecker_support} in
    \doclink{TNLean/Channel/Schwarz/PetzEqualityPrerequisites}
    (\path{TNLean/Channel/Schwarz/PetzEqualityPrerequisites.lean}). It assumes
    positive semidefiniteness, $\ker\sigma\subseteq\ker\rho$, and unit trace of
    $\tau$, but does not require $\rho$ or $\sigma$ to be normalized. Its proof
    uses the tensor-logarithm formula on support projections and support
    absorption. These ancilla results use only the standard \Lean{} axioms and
    no Lieb input.

    The twirling step appends the maximally mixed state
    $\tau=\Id_C/d_C$ and represents the partial trace by a finite unitary
    $1$-design. Because the current dependency supplies no Haar measure on the
    matrix unitary group, the formalization uses the Heisenberg--Weyl operators
    $W(a,b)=X^aZ^b$. For a primitive $d$-th root of unity,
    \leanid{Matrix.sum_weyl_conj} in
    \doclink{TNLean/Channel/Schwarz/WeylTwirl}
    (\path{TNLean/Channel/Schwarz/WeylTwirl.lean}) proves
    \begin{align}
        \frac{1}{d^2}\sum_{a,b}W(a,b)MW(a,b)^\dagger
            &= \frac{\tr M}{d}\Id.
        \label{eq:cpsv16_ssa_from_lieb_route_weyl_twirl}
    \end{align}
    The proof uses root-of-unity character-sum orthogonality,
    \leanid{Matrix.sum_rootOfUnity_pow_eq}, and no Lieb input.

    The ancilla-factor lift
    \leanid{Matrix.sum_kronecker_one_weyl_conj} in
    \doclink{TNLean/Channel/Schwarz/RelativeEntropyDataProcessing}
    (\path{TNLean/Channel/Schwarz/RelativeEntropyDataProcessing.lean}) proves
    \begin{align}
        \frac{1}{d_C^2}\sum_{a,b}
        (\Id_S\otimes W(a,b))M(\Id_S\otimes W(a,b))^\dagger
            &= (\tr_C M)\otimes\frac{\Id_C}{d_C}.
        \label{eq:cpsv16_ssa_from_lieb_route_partial_trace_twirl}
    \end{align}
    Joint convexity applied to
    Eq.~\ref{eq:cpsv16_ssa_from_lieb_route_partial_trace_twirl}, together with
    Eqs.~\ref{eq:cpsv16_ssa_from_lieb_route_unitary_invariance} and
    \ref{eq:cpsv16_ssa_from_lieb_route_ancilla_additivity}, gives
    \begin{align}
        D(\tr_C\rho\,\| \tr_C\sigma)
            &\leq D(\rho\,\| \sigma).
        \label{eq:cpsv16_ssa_from_lieb_route_partial_trace_monotonicity}
    \end{align}
    The positive-definite theorem is
    \leanid{quantumRelativeEntropy_partialTraceRight_le} in
    \doclink{TNLean/Channel/Schwarz/RelativeEntropyDataProcessing}. Transport
    of joint convexity to the product index uses
    \leanid{quantumRelativeEntropy_submatrix_equiv}.

    The support-domain theorem is
    \leanid{quantumRelativeEntropy_partialTraceRight_le_support} in the same
    file. It regularizes both arguments by
    \begin{align}
        M_\varepsilon
            &= \frac{M+\varepsilon\Id}{1+N_{SC}\varepsilon},
        \qquad \varepsilon>0.
        \label{eq:cpsv16_ssa_from_lieb_route_data_processing_regularization}
    \end{align}
    The positive-definite theorem applies for every $\varepsilon>0$. The
    right-hand side converges to $D(\rho\,\| \sigma)$ by the two-argument
    limit. The partial traces are affine regularizations of
    $\tr_C\rho$ and $\tr_C\sigma$ with scaling rate $N_{SC}$ and shift rate
    $d_C$; this identity is
    \leanid{partialTraceRight_regPerturbAffine}. The support condition passes
    to the marginals by \leanid{Matrix.partialTraceRight_support}. The
    arbitrary-rate affine limit
    \leanid{TNLean.RelativeEntropyConvexity.tendsto_relativeEntropyPerturbAffine}
    applies to maps of the form
    \begin{align}
        M
            &\longmapsto
            \frac{M+b\varepsilon\Id}{1+a\varepsilon},
        \qquad a,b>0.
        \label{eq:cpsv16_ssa_from_lieb_route_arbitrary_affine_regularization}
    \end{align}
    Passing the inequality through both limits proves
    Eq.~\ref{eq:cpsv16_ssa_from_lieb_route_partial_trace_monotonicity} whenever
    $\ker\sigma\subseteq\ker\rho$. This extension uses no Lieb input beyond the
    positive-definite joint-convexity theorem.

    \item \emph{Final assembly.}

    Define
    \begin{align}
        \sigma_{ABC}
            &= \frac{\Id_A}{d_A}\otimes\rho_{BC},
        \qquad d_A=\dim\mathcal H_A.
        \label{eq:cpsv16_ssa_from_lieb_route_reference_state}
    \end{align}
    If $\rho_{BC}$ is singular, normalization alone does not put
    $(\rho_{ABC},\sigma_{ABC})$ in the domain of relative entropy. The required
    condition is $\ker\sigma_{ABC}\subseteq\ker\rho_{ABC}$.

    The marginal support lemma states
    \begin{align}
        \mathcal H_A\otimes\ker\rho_{BC}
            &\subseteq\ker\rho_{ABC},
        \label{eq:cpsv16_ssa_from_lieb_route_marginal_kernel}\\
        \operatorname{supp}\rho_{ABC}
            &\subseteq
            \mathcal H_A\otimes\operatorname{supp}\rho_{BC}.
        \label{eq:cpsv16_ssa_from_lieb_route_marginal_support}
    \end{align}
    Let $P_{BC}$ be the projection onto
    $\operatorname{supp}\rho_{BC}$, and define
    \begin{align}
        Q
            &= \Id_A\otimes(\Id_{BC}-P_{BC}).
        \label{eq:cpsv16_ssa_from_lieb_route_kernel_projection}
    \end{align}
    Since $\tr_A\rho_{ABC}=\rho_{BC}$ and
    $(\Id_{BC}-P_{BC})\rho_{BC}=0$, the partial-trace adjoint identity gives
    \begin{align}
        \tr(Q\rho_{ABC})
            &= \tr((\Id_{BC}-P_{BC})\tr_A\rho_{ABC})
        \label{eq:cpsv16_ssa_from_lieb_route_projection_trace_reduction}\\
            &= \tr((\Id_{BC}-P_{BC})\rho_{BC})
        \label{eq:cpsv16_ssa_from_lieb_route_projection_trace_marginal}\\
            &= 0.
        \label{eq:cpsv16_ssa_from_lieb_route_projection_trace_zero}
    \end{align}
    Positivity of $\rho_{ABC}$ and idempotence of $Q$ imply
    \begin{align}
        \tr(Q\rho_{ABC})
            &= \lVert Q\rho_{ABC}^{1/2}\rVert_2^2,
        \label{eq:cpsv16_ssa_from_lieb_route_projection_hilbert_schmidt}\\
        Q\rho_{ABC}^{1/2}
            &= 0,
        \label{eq:cpsv16_ssa_from_lieb_route_projection_square_root_zero}\\
        Q\rho_{ABC}
            &= 0.
        \label{eq:cpsv16_ssa_from_lieb_route_projection_state_zero}
    \end{align}
    Thus the range of $Q$, namely
    $\mathcal H_A\otimes\ker\rho_{BC}$, lies in $\ker\rho_{ABC}$, proving
    Eqs.~\ref{eq:cpsv16_ssa_from_lieb_route_marginal_kernel} and
    \ref{eq:cpsv16_ssa_from_lieb_route_marginal_support}. Because
    $\Id_A/d_A$ is full rank,
    \begin{align}
        \ker\sigma_{ABC}
            &= \mathcal H_A\otimes\ker\rho_{BC}
            \subseteq\ker\rho_{ABC}.
        \label{eq:cpsv16_ssa_from_lieb_route_reference_kernel_condition}
    \end{align}

    The marginal support lemma is
    \leanid{Matrix.marginal_support} in
    \doclink{TNLean/Analysis/MarginalSupport}
    (\path{TNLean/Analysis/MarginalSupport.lean}). It states that if a
    Hermitian idempotent on $B\otimes C$ annihilates
    $\tr_A\rho_{ABC}$ for positive semidefinite $\rho_{ABC}$, then its
    identity-on-$A$ lift annihilates $\rho_{ABC}$. Its operator core is
    \leanid{Matrix.PosSemidef.proj_mul_eq_zero_of_trace_eq_zero}, and its
    partial-trace adjoint identity is
    \leanid{Matrix.traceA_ABC_lift_trace}. These results use only the standard
    \Lean{} axioms.

    On $\operatorname{supp}\sigma_{ABC}$, the logarithm splits as
    \begin{align}
        \log\left(\frac{\Id_A}{d_A}\otimes\rho_{BC}\right)
            &= -(\log d_A)\Id_{ABC}
              +\Id_A\otimes\log\rho_{BC}.
        \label{eq:cpsv16_ssa_from_lieb_route_reference_logarithm}
    \end{align}
    Here $\log\rho_{BC}$ is restricted to
    $\operatorname{supp}\rho_{BC}$. The support condition ensures that
    $\rho_{ABC}$ has no weight outside this subspace. Therefore,
    \begin{align}
        D\left(\rho_{ABC} \middle\|
          \frac{\Id_A}{d_A}\otimes\rho_{BC}\right)
            &= \log d_A+S(\rho_{BC})-S(\rho_{ABC}).
        \label{eq:cpsv16_ssa_from_lieb_route_tripartite_relative_entropy}
    \end{align}

    Under $\Lambda=\tr_C$, the two arguments become
    \begin{align}
        \Lambda(\rho_{ABC})
            &= \rho_{AB},
        \label{eq:cpsv16_ssa_from_lieb_route_partial_trace_state}\\
        \Lambda(\sigma_{ABC})
            &= \frac{\Id_A}{d_A}\otimes\rho_B.
        \label{eq:cpsv16_ssa_from_lieb_route_partial_trace_reference}
    \end{align}
    Applying the marginal support lemma to $\rho_{AB}$ and
    $\rho_B=\tr_A\rho_{AB}$ puts the image pair in the same relative-entropy
    domain. Its relative entropy is
    \begin{align}
        D\left(\rho_{AB} \middle\|
          \frac{\Id_A}{d_A}\otimes\rho_B\right)
            &= \log d_A+S(\rho_B)-S(\rho_{AB}).
        \label{eq:cpsv16_ssa_from_lieb_route_bipartite_relative_entropy}
    \end{align}
    Substituting
    Eqs.~\ref{eq:cpsv16_ssa_from_lieb_route_tripartite_relative_entropy} and
    \ref{eq:cpsv16_ssa_from_lieb_route_bipartite_relative_entropy} into
    partial-trace monotonicity gives
    \begin{align}
        \log d_A+S(\rho_B)-S(\rho_{AB})
            &\leq
            \log d_A+S(\rho_{BC})-S(\rho_{ABC}),
        \label{eq:cpsv16_ssa_from_lieb_route_entropy_data_processing}\\
        S(\rho_B)-S(\rho_{AB})
            &\leq S(\rho_{BC})-S(\rho_{ABC}),
        \label{eq:cpsv16_ssa_from_lieb_route_conditional_entropy_order}\\
        S(\rho_{ABC})+S(\rho_B)
            &\leq S(\rho_{AB})+S(\rho_{BC}).
        \label{eq:cpsv16_ssa_from_lieb_route_ssa_assembly}
    \end{align}
    The final line is
    Eq.~\ref{eq:cpsv16_ssa_from_lieb_route_strong_subadditivity}, equivalently
    the nonnegativity of the conditional mutual information
    $I(A:C\mid B)\geq0$.

    This construction discards $C$ and places the maximally mixed state on the
    retained factor $A$. Discarding $A$ against
    $\rho_A\otimes\rho_{BC}$ would prove only subadditivity,
    $I(A:BC)\geq0$, rather than strong subadditivity.
\end{enumerate}

\section{Formal completion of the route}

Layers (1)--(6) are formalized. The route begins with the entropy bridge,
defines \leanid{quantumRelativeEntropy}, proves Klein nonnegativity on the full
support domain, derives joint convexity from Lieb concavity, establishes
partial-trace data processing, and instantiates that result at the reference
state in Eq.~\ref{eq:cpsv16_ssa_from_lieb_route_reference_state}.

The positive-definite final theorem is
\leanid{strong_subadditivity_posDef} in
\doclink{TNLean/Channel/Schwarz/StrongSubadditivityPosDef}
(\path{TNLean/Channel/Schwarz/StrongSubadditivityPosDef.lean}). In this case,
$\rho_{BC}$ and the reference state
$(\Id_A/d_A)\otimes\rho_{BC}$ are positive definite, so the marginal support
lemma is unnecessary. The tensor-logarithm identity and the partial-trace
adjoint relation evaluate both relative entropies, and the positive-definite
data-processing theorem yields strong subadditivity.

The general theorem is \leanid{strong_subadditivity_general} in the same file.
It applies to every positive semidefinite unit-trace $\rho_{ABC}$. The theorem
\leanid{Matrix.marginal_support} places the possibly singular reference state
in the domain of
\leanid{quantumRelativeEntropy_partialTraceRight_le_support}. The two relative
entropies are evaluated by \leanid{rel_entropy_eval_support}, which regularizes
the reduced state as
\begin{align}
    \rho_{R,\varepsilon}
        &= \frac{\rho_R+\varepsilon\Id}{1+d_R\varepsilon},
        \qquad \varepsilon>0,
    \label{eq:cpsv16_ssa_from_lieb_route_reduced_state_regularization}
\end{align}
and passes to the limit. This limit argument is required because the direct
tensor-logarithm split is unavailable for a singular reference state.

The stable entropy interface re-exports the inequality from
\doclink{TNLean/Entropy/StrongSubadditivity}. Its consumers import that theorem
rather than the former axiom module. Their statements require no change; only
the justification has changed from the axiom \leanid{strong_subadditivity} to
the theorem \leanid{strong_subadditivity_general}.

\section{Verdict}

Strong subadditivity is proved on the full density-matrix domain. The route is
the classical relative-entropy argument: the entropy bridge, quantum relative
entropy, Klein's inequality, joint convexity from Lieb concavity, data
processing under a partial trace, and the tripartite instantiation.

The positive-definite joint-convexity theorem
\leanid{convexOn_quantumRelativeEntropy} is extended to the support domain by
\leanid{convexOn_quantumRelativeEntropy_support}. Unitary invariance,
ancilla-additivity, the finite Heisenberg--Weyl twirl, and its ancilla lift give
\leanid{quantumRelativeEntropy_partialTraceRight_le}; regularization gives
\leanid{quantumRelativeEntropy_partialTraceRight_le_support}. The marginal
support lemma and \leanid{rel_entropy_eval_support} then establish
\leanid{strong_subadditivity_general} for every positive semidefinite
unit-trace state.

The only entropy-side dependency of this route that replaces the former axiom
is the proved theorem \leanid{lieb_concavity_posDef}. The theorem
\leanid{strong_subadditivity_general} has exactly the hypotheses of the former
standalone axiom. Consequently, that axiom has been removed, and
\leanid{Entropy.strongSubadditivity} now re-exports the derived theorem.

\begin{thebibliography}{9}

\bibitem{LiebRuskai1973}
  E.~H.~Lieb and M.~B.~Ruskai,
  \emph{Proof of the strong subadditivity of quantum-mechanical entropy},
  Journal of Mathematical Physics \textbf{14} (1973), 1938--1941.

\bibitem{Lieb1973Convex}
  E.~H.~Lieb,
  \emph{Convex trace functions and the Wigner--Yanase--Dyson conjecture},
  Advances in Mathematics \textbf{11} (1973), 267--288.

\end{thebibliography}

\end{document}
